\section{Introduction}

This project investigates a pragmatic question: can a sequence model trained offline on nonlinear trajectory-optimization solutions provide warm-starts that reduce the cost of solving a minimum-time raceline problem with obstacle avoidance? The project focus is narrower than a generic reinforcement learning study. The relevant downstream objective is not imitation loss alone, but whether a DT-generated initial trajectory helps the nonlinear optimizer solve faster, with fewer iterations, or with higher robustness than the baseline initializer. The setup combines nonlinear vehicle trajectory optimization with Decision Transformer-style return-conditioned sequence modeling \cite{chen2021decision,subositsgerdes2021fidelity}.

The final repo state reflects a progressively narrowed scope. The active execution loop is Oval-track only, uses FATROP-clean generated data for training shards, and uses fixed start-state assumptions for evaluation: the start progress $s_0$ varies along the centerline, while $e_0=d\psi_0=u_{y,0}=r_0=0$ and $u_{x,0}=5.0$ m/s. Two benchmark families are used: a no-obstacle family and a frozen 1--4 obstacle family. In practice, IPOPT was used for most completed diagnostic evaluations because full fixed-gate FATROP evaluation was too slow for interactive iteration, while FATROP remained the generation path for the cleaned training data.

The main result is negative but informative. The implemented DT system learns the offline data well and successive data interventions do improve validation loss. However, those gains do not produce a convincing end-to-end warm-start benefit. This report therefore emphasizes what was actually implemented, what was measured, which confounds were removed, and what the final failure mode appears to be.
